\section{事件时间仿真与研究项目}
\label{ch:lower-microstructure-simulation}

\begin{itemize}
  % Generated from curriculum/manifest.json; do not edit by hand.
\item 直接先修：下册《执行与做市：价格冲击与库存反馈》。

\end{itemize}

仿真将到达、报价、成交和库存连续串起来。比较策略时要重复生成市场路径，并保留同一路径上的结果差，才能区分策略作用与随机波动。


\MFQTransition{先确定生成机制，再比较策略}

% Generated by tools/build_figures.py; edit figures/figure-specs.json instead.
\MFQFigure{tex/figures/generated/lower-ms-execution.pdf}{执行轨迹在冲击成本与库存风险之间权衡；前置成交降低尾部库存，却提高早期冲击。}{lower-ms-execution}


一个最小事件生成器需要明确：等待时间、订单方向、价格变化、委托大小和终止条件。例如
\[
 \Delta t_i\sim\operatorname{Exp}(\lambda),\quad
 q_i\in\{-1,+1\},\quad
 \Delta m_i=\beta q_i+\sigma\sqrt{\Delta t_i}\varepsilon_i.
\]
这些方程是数据生成假设，不是对现实的事实陈述。若要引入季节或 Hawkes 聚集，应逐层增加并分别做残差诊断；模型复杂度不能替代可识别性。

比较静态报价与库存控制时，可以让两者使用相同的 $(\Delta t_i,q_i,\varepsilon_i)$ 和成交均匀数，但不能强制共享最终成交结果。每个策略用自己的报价计算成交阈值，再把同一个均匀数与阈值比较；因此报价能够改变成交、库存与收益，同时随机源仍可配对。共享底层随机数有助于减少比较的噪声；独立路径也能形成有效估计，只是需要正确计算两组的不确定性。共享随机源不等于强制两种策略具有相同成交。

\MFQTransition{共同随机数何时降低方差}

对同一路径上的结果 $Y_A,Y_B$，差值 $D=Y_A-Y_B$ 满足
\[
 \operatorname{Var}D=\operatorname{Var}Y_A+\operatorname{Var}Y_B
 -2\operatorname{Cov}(Y_A,Y_B).
\]
若方差都为 4、相关系数为 0.75，则差值方差为 $8-6=2$；独立抽样则为 8。100 条独立配对路径的均值差标准误为 $\sqrt{2/100}$，是独立比较的一半。若协方差为负，配对反而可能增加方差。

一个成交例子也能说明为什么不应共享最终成交：同一均匀数为 0.4，A 报价对应成交概率 0.3，B 对应 0.6，则 A 不成交、B 成交。底层随机源相同，策略后果不同。notebook 用这一规则比较静态与库存偏移报价，在多条独立路径上计算损益差及最大库存，而不把一条路径的胜负当作总体结论。

做市商的损益由现金与期末库存估值构成：主动买单在卖价成交时，做市商收到现金、库存减少；主动卖单在买价成交时方向相反。库存反馈规则先根据当前库存调整报价，报价又影响成交概率，成交后的库存再进入下一次决策。配套实验在多条路径上比较这些递推产生的损益与库存分布。

\MFQTransition{同时比较损益与库存风险}

把终点累计损益记为 $\operatorname{PnL}_T$。可同时考察
\[
 \text{PnL},\quad \max_t|q_t|,\quad q_T,\quad
 \text{fill rate},\quad \text{turnover},\quad \text{shortfall},
\]
这些量的路径分布描述平均表现以外的风险。若只看平均 PnL，策略可以通过无限库存或未完成母单获得虚假优势。做市 PnL 还应拆成点差捕获、持仓重估、费用与冲击；执行策略则应以相同到达价和母单任务比较。

动态规划策略与仿真统计承担不同证据：DP 在给定离散状态、转移与目标下给出最优控制；Monte Carlo 描述该控制在声明分布下的结果分布。DP 最优不意味着模型正确，仿真均值领先也不意味着总体或未来市场领先。

\MFQTransition{公开汇总数据能说明什么}

配套材料使用公开的 SEC 汇总统计设定一个低成交率的比较情形。把其数值作为模型输入后，实际模拟成交仍由每种策略自己的报价与成交规则决定。

公开汇总数据可以帮助设定低成交率等比较情形，但没有提供具体订单的队列位置、交易发起方或网络延迟。由撤单成交比率换算为执行概率，还额外假设每个终止事件只有取消或成交两类结果。因此，这种换算只能在所给简化模型内解释。

\MFQTransition{比较对象要一致}
两种策略从相同初始资金和库存出发，各自按报价决定成交。期末库存按同一价格估值，费用与未完成任务采用相同口径。这样才能把损益差解释为策略在所给模型中的差异\cite{hasbrouck2007empirical}。

\MFQTransition{事件先后顺序与策略可用信息}

事件时间模型按发生顺序推进：市场变化形成新的可见信息，策略据此报价或下单，订单经过延迟后才可能成交。把数据到达、决策、交易所确认和成交分开，才能表达动作发生后、成交之前的价格风险。

外生订单流、价格创新、延迟和成交抽样可以分别生成。使用共同随机数比较策略时，共享外生随机变化仍允许各策略采取不同动作；由此产生的成交和库存路径可以分叉。

\MFQTransition{校准与验证不是同一件事}

校准让模拟器匹配选定统计量，例如到达率、点差、深度、订单大小和自相关；验证检查未参与校准的统计量，例如价格响应、队列等待和极端流动性。若所有统计量都用于拟合，就没有留出证据。

一条仿真路径同时受到数量与时间关系约束：成交数量不能超过可成交委托，现金变化应与成交方向一致，未来事件不能影响过去的决策。均值、尾部和时间依赖则描述多条路径的分布性质。个别路径上的数量关系与整体分布是否合适，需要分别考察。

\MFQTransition{策略比较与置信区间}

对策略 A/B 使用相同外生随机路径，比较配对差 $D_b=Y_b^A-Y_b^B$，若两个结果正相关，其标准误可小于两组独立模拟。报告均值差、置信区间、尾部分位和失败率。若策略改变市场状态，后续事件可能分叉，但仍可共享底层随机数。

调参路径与最终评价路径分开。增加仿真次数只降低给定模拟器内的 Monte Carlo 误差，不降低模拟器错设。结论应写成“在这些生成假设下”，而不是“市场中”。

\MFQTransition{压力场景与代理交互}

压力场景包括到达率骤降、撤单激增、价格跳跃、延迟上升、单边流动性枯竭和交易所拒单。这些情形分别改变流动性、价格或信息延迟；比较它们有助于找出策略依赖最强的假设。

多代理仿真可加入做市者、噪声交易者、执行者和信息交易者。产生相似 stylized facts 不代表机制被识别：不同代理规则可能生成相同分布。应做消融，逐类移除代理并观察哪些统计量变化。

\MFQTransition{历史回放与反事实的局限}

历史消息回放可以重现已经发生的事件顺序；实时观察而不实际下单，可以检验信号能否在相应时点形成。两者都没有观察到假想订单真正提交后会产生的成交与市场影响，因此这部分仍依赖模型假设。

把仿真结果推广到市场，需要考虑订单对流动性和其他参与者行为的反馈。本书的综合项目保留这些限制，重点比较给定模型内的策略后果。

\begin{diagnostic}
某条路径上的领先可能来自随机波动。增加独立路径可以估计平均差异；更换流动性状态或生成机制，则用来研究结论对模型假设的依赖。两种比较回答不同的问题。
\end{diagnostic}

\section{本路线的综合项目}

选择一个报价或执行问题，把事件流、成交和库存递推连起来。用多条路径比较策略，给出损益差的置信区间，并观察延迟及未成交怎样改变结果。再说明哪些假设能够由历史消息检验，哪些涉及尚未观察到的订单后果。

\section{综合练习}
\begin{enumerate}[leftmargin=*]
  \item \textbf{口述概念：}条件强度与无条件平均事件率有什么差别？
  \item \textbf{笔试推导：}由指数等待时间似然推导 $\widehat\lambda=n/\sum_i\Delta_i$。
  \item \textbf{笔试推导：}解释 Hawkes 似然中补偿子为何不可省略。
  \item \textbf{数值编程：}扫描 $\alpha/\beta$ 并画出单个事件后的强度衰减。
  \item \textbf{订单簿：}手工推进两档 FIFO 簿，记录撤单、部分成交和未成交。
  \item \textbf{研究判断：}找出把 maker side 当主动方时报告会发生的错误。
  \item \textbf{错误研究材料：}某报告只用成交笔数拟合 Hawkes，却不处理盘中季节性；列出至少三项缺失证据。
  \item \textbf{开放任务：}设计一个带盘口深度、撤单和队列位置的数据协议，明确可知时点和失败状态。
\end{enumerate}

\MFQLead{执行、做市与库存控制题组}
\begin{enumerate}[leftmargin=*]
  \item \textbf{口述概念：}临时冲击与永久冲击怎样进入不同期的成交价？
  \item \textbf{笔试推导：}手算冻结日程 $(3,2,0)$ 的 partial-fill IS。
  \item \textbf{笔试推导：}写出 DP 终端条件，并解释为何剩余库存不能零成本消失。
  \item \textbf{数值编程：}枚举六手三期所有日程，与 DP 结果逐项比较。
  \item \textbf{控制实验：}提高库存惩罚，记录报价中心与期末库存如何变化。
  \item \textbf{研究判断：}识别“先生成报价、后计算库存”的前视错误。
  \item \textbf{错误研究包：}某执行报告只给平均成交价，没有到达价、未完成量和停止原因；写出修复清单。
  \item \textbf{开放任务：}加入随机完成率与最大损失停止，设计可复现的压力矩阵。
\end{enumerate}

\MFQTransition{事件仿真、策略比较与综合项目题组}
\begin{enumerate}[leftmargin=*]
  \item \textbf{口述概念：}为什么共同随机数能提高控制器差异比较的精度？
  \item \textbf{笔试推导：}写出静态半点差策略在给定方向与价格变化下的逐事件 PnL。
  \item \textbf{统计设计：}为 PnL、期末库存和最大库存设计置信区间。
  \item \textbf{数值编程：}把齐次 Poisson 换为日内季节强度，并验证时间变换残差。
  \item \textbf{反例实验：}重复记录一次成交或改变事件先后顺序，手算库存与现金会怎样出错，并指出被破坏的数量或信息关系。
  \item \textbf{研究判断：}审阅只报告单路径 PnL 的做市研究材料。
  \item \textbf{真实数据：}解释 SEC 聚合 cancel-to-trade 比率为何不能识别自己的队列成交概率。
  \item \textbf{综合项目：}建立报价、成交与库存相互影响的小模型，比较两种策略，并解释参数变化怎样影响结论。
\end{enumerate}

\subsection*{基础衔接：独立计算}
两策略在同一路径上的损益完全相同，配对差标准误为零能否证明市场盈利？

